\chapter{Appendix A: Muscle Force and Joint Torque Calculations}\label{app:muscle_calcs}
\graphicspath{{../Figures/92-AppendixA/}{figs/92-AppendixA/}}

Chapter~\ref{ch:torque} uses isometric muscle force to define human knee-torque targets. In Chapter~\ref{ch:sns}, the model must also determine how force changes as a muscle shortens or lengthens under neural control. This appendix describes the human muscle formulation in Gait2392, the tension law used after conversion to MuJoCo, and the separate pressure-dependent force and compliance calculations used for the robot actuators in Chapters~\ref{ch:muscles} and~\ref{ch:torque}.

\section{Human Muscle Force Calculations}
\subsection{OpenSim Thelen2003 and Millard2012 Formulations}\label{sec:app_opensim_curves}

The OpenSim Gait2392 human model, its subject01 version used for static optimization, and the rerouted robot counterpart each specify \texttt{Thelen2003Muscle} for all 92 muscles. The Millard equilibrium muscle provides a comparison with another way of representing active fibers, passive tissue, and a tendon in series \citep{thelen_adjustment_2003,millard_flexing_2013}. For either formulation, the force and length balances are
\begin{equation}
 F_T(l_T)=F_M(l_M,\dot l_M,a)\cos\alpha,
 \qquad l_{MT}=l_T+l_M\cos\alpha,
 \label{eq:app_muscle_equilibrium}
\end{equation}
where $a$ is activation, $\alpha$ is pennation, and $l_M$, $l_T$, and $l_{MT}$ are fiber, tendon, and musculotendon lengths. The cosine projects fiber force onto the tendon direction. When tendon compliance is enabled, fiber length is a state determined by the equilibrium calculation rather than the full musculotendon path length.

In the damped Millard formulation, fiber force has the form
\begin{equation}
 F_M=F_0\left[a f_L(\tilde l_M)f_V(\tilde v_M)
 + f_{PE}(\tilde l_M)+\beta\tilde v_M\right],
 \label{eq:app_millard_force}
\end{equation}
where $F_0$ is maximum isometric fiber force, $\tilde l_M$ is fiber length divided by optimal length, and $\tilde v_M$ is fiber velocity divided by maximum shortening speed. The dimensionless coefficient $\beta$ scales the velocity-dependent damping force. The next two sections give the characteristic curves that turn these quantities into active and passive tension. A forward simulation can advance activation and fiber length in time. The static-optimization forces used in Section~\ref{sec:muscle-models-compare} were estimated separately at each sampled pose \citep{opensim_static_optimization_docs}.

\subsubsection{Thelen activation and contraction dynamics}

A change in neural excitation first changes activation. Fiber length and shortening or lengthening speed then determine how much tension that activation can produce \citep{thelen_adjustment_2003}. The equations below follow the Thelen formulation with the tensile branches used in OpenSim \citep{opensim_muscle_curves_44}. A bar denotes normalization by maximum isometric fiber force $F_0^M$ or by optimal fiber length $L_0^M$, as appropriate.

Activation follows
\begin{equation}
 \frac{da}{dt}=\frac{u-a}{\tau_a(a,u)},
 \label{eq:app_thelen1}
\end{equation}
with an activation-dependent time constant
\begin{equation}
 \tau_a(a,u)=\begin{cases}
 \tau_{\mathrm{act}}(0.5+1.5a),&u>a,\\
 \tau_{\mathrm{deact}}/(0.5+1.5a),&u\leq a.
 \end{cases}
 \label{eq:app_thelen2}
\end{equation}
Here $u\in[0,1]$ is neural excitation, $a$ is muscle activation, and $\tau_{\mathrm{act}}$ and $\tau_{\mathrm{deact}}$ set the activation and deactivation time scales.

The normalized passive fiber force and active force--length multiplier are
\begin{equation}
 \bar F^{PE}=\begin{cases}
 0,&\bar L^M\leq1,\\[2pt]
 \dfrac{\exp\!\left[k^{PE}(\bar L^M-1)/\varepsilon_0^M\right]-1}
 {\exp(k^{PE})-1},&\bar L^M>1,
 \end{cases}
 \label{eq:app_thelen3}
\end{equation}
\begin{equation}
 f_l=\exp\!\left[-\frac{(\bar L^M-1)^2}{\gamma}\right].
 \label{eq:app_thelen4}
\end{equation}
The tendon relation is
\begin{equation}
 \bar F^T=\begin{cases}
 0,&\varepsilon^T\leq0,\\[2pt]
 \displaystyle\frac{\bar F^T_{\mathrm{toe}}}{\exp(k_{\mathrm{toe}})-1}
 \left[\exp\!\left(\frac{k_{\mathrm{toe}}\varepsilon^T}{\varepsilon^T_{\mathrm{toe}}}\right)-1\right],
 &0<\varepsilon^T\leq\varepsilon^T_{\mathrm{toe}},\\[6pt]
 k_{\mathrm{lin}}(\varepsilon^T-\varepsilon^T_{\mathrm{toe}})+\bar F^T_{\mathrm{toe}},
 &\varepsilon^T>\varepsilon^T_{\mathrm{toe}}.
 \end{cases}
 \label{eq:app_thelen5}
\end{equation}
For a given tendon force, the required normalized active fiber force is $\bar F^M=\bar F^T/\cos\alpha-\bar F^{PE}$. Thelen expresses the corresponding fiber velocity as
\begin{equation}
 V^M=(0.25+0.75a)V^M_{\max}\frac{\bar F^M-a f_l}{b},
 \label{eq:app_thelen6}
\end{equation}
\begin{equation}
 b=\begin{cases}
 a f_l+\bar F^M/A_f,&\bar F^M\leq a f_l,\\[4pt]
 \displaystyle\frac{(2+2/A_f)(a f_l\bar F^M_{\mathrm{len}}-\bar F^M)}
 {\bar F^M_{\mathrm{len}}-1},&\bar F^M>a f_l.
 \end{cases}
 \label{eq:app_thelen7}
\end{equation}

Here $\bar F^M$ is the normalized active fiber force entering the force--velocity relation, $V^M=\dot L^M$ is negative during shortening, and $V^M_{\max}$ is the maximum shortening-speed scale. Table~\ref{tab:app_thelen_symbols} defines the curve parameters. A shape factor is a dimensionless parameter controlling a curve's curvature or width; it is not a stiffness in \unit{\newton\per\meter}. For example, increasing $\gamma$ broadens the active force--length curve while preserving its peak at $\bar L^M=1$. Changing $k^{PE}$ changes the passive curve's shape while retaining unit force at strain $\varepsilon_0^M$.

\begin{table}[htbp]
\centering\small
\caption{Length, strain, and shape parameters in the Thelen formulation.}
\label{tab:app_thelen_symbols}
\begin{tabularx}{\textwidth}{@{}lX@{}}
\toprule
Symbol & Definition \\
\midrule
$\bar L^M=L^M/L_0^M$ & Fiber length relative to optimal fiber length. \\
$\varepsilon_0^M$ & Passive fiber strain at force $F_0^M$. \\
$k^{PE}$ & Passive force--length exponential shape factor (5 in the paper). \\
$\gamma$ & Active force--length width parameter (0.45 in the paper). \\
$\varepsilon^T$ & Tendon strain relative to tendon slack length. \\
$\varepsilon_0^T$ & Tendon strain at force $F_0^M$. \\
$\varepsilon^T_{\mathrm{toe}}$ & Strain at the end of the nonlinear tendon region, approximately $0.609\varepsilon_0^T$. \\
$\bar F^T_{\mathrm{toe}}$ & Normalized force at that transition; 0.33. \\
$k_{\mathrm{toe}},\ k_{\mathrm{lin}}$ & Tendon exponential shape factor (3) and linear slope ($1.712/\varepsilon_0^T$). \\
$A_f$ & Force--velocity shape factor (0.25 in the paper). \\
$\bar F^M_{\mathrm{len}}$ & Maximum normalized lengthening-force factor. \\
\bottomrule
\end{tabularx}
\end{table}

The force--velocity multiplier also follows directly from Equations~\eqref{eq:app_thelen6} and~\eqref{eq:app_thelen7}. Let $w=V^M/[(0.25+0.75a)V^M_{\max}]$ and $c=(2+2/A_f)/(\bar F^M_{\mathrm{len}}-1)$. On the shortening and lengthening branches, respectively,
\begin{equation}
 f_V=\frac{\bar F^M}{a f_l}=\begin{cases}
 (1+w)/(1-w/A_f),&-1\leq w\leq0,\\
 (1+c\bar F^M_{\mathrm{len}}w)/(1+cw),&w>0.
 \end{cases}
 \label{eq:app_thelen_velocity_multiplier}
\end{equation}
Both branches give $f_V=1$ at zero velocity. Shortening reduces tension, whereas lengthening increases it toward $\bar F^M_{\mathrm{len}}$. The OpenSim implementation uses linear continuation outside its regular force range so that the velocity relation remains invertible \citep{opensim_muscle_curves_44}. The Gait2392 files specify $\gamma=0.5$, $k^{PE}=4$, $A_f=0.3$, $\bar F^M_{\mathrm{len}}=1.8$, $\varepsilon_0^M=0.6$, and $\varepsilon_0^T=0.033$. For these parameters the tendon carries $F_0^M$ at 3.3\,\% elongation. Table~\ref{tab:app_thelen_symbols} also identifies the values used in the original publication.

\subsubsection{Millard equilibrium and damping}

In the model of \citet{millard_flexing_2013}, the fiber velocity must satisfy the series force balance
\begin{equation}
 \big[a f_L(x)f_V(v)+f_{PE}(x)+\beta v\big]\cos\alpha=f_T(y),
 \qquad x=\frac{l_M}{l_{\mathrm{opt}}},\quad
 y=\frac{l_T}{l_{\mathrm{slack}}},\quad v=\frac{\dot l_M}{v_{\mathrm{max}}^M}.
 \label{eq:app_millard_equilibrium}
\end{equation}
Here $v_{\mathrm{max}}^M$ is the maximum shortening speed in units of length per second, and $v<0$ denotes shortening. Activation scales the active fiber force. Passive fibers can develop tension at zero activation, and the tendon stretches until its tension equals the component of fiber force along the tendon. The damping term, $\beta v$, makes fiber velocity affect equilibrium even at low activation. The undamped formulation sets $\beta=0$.

OpenSim represents each of the four characteristic curves with joined quintic B\'ezier segments \citep{millard_flexing_2013,opensim_muscle_curves_44}. On one segment the dimensionless input $s$ and normalized force $f$ are
\begin{equation}
 (s(u),f(u))=\sum_{i=0}^{5}\binom{5}{i}(1-u)^{5-i}u^i(X_i,Y_i),
 \qquad 0\leq u\leq1.
 \label{eq:app_millard_bezier}
\end{equation}
The six control points $(X_i,Y_i)$ on each segment are set by the curve's lengths or speeds, force levels, slopes, and curvature settings. To evaluate a curve at a given length or speed, OpenSim locates its segment, solves $s(u)$ for $u$, and evaluates $f(u)$. The segments meet with continuous first and second derivatives. Table~\ref{tab:app_millard_curves} gives the values from the default damped equilibrium muscle in OpenSim~4.4. Its dimensionless damping coefficient is $\beta=0.1$.

\begin{table}[htbp]
\centering\small
\caption{Characteristic curves of the default damped Millard muscle. Force is normalized by maximum isometric fiber force. Shortening velocity is negative.}
\label{tab:app_millard_curves}
\begin{tabularx}{\textwidth}{@{}p{0.23\textwidth}X@{}}
\toprule
Curve & Values and mechanical interpretation \\
\midrule
Active fiber, $f_L(x)$ & Peaks at $f_L(1)=1$. At $x=0.75$ and $1.20$, the multipliers are 0.826 and 0.824. The curve reaches zero at the short and long shoulders, $x=0.4441$ and $1.8123$. \\
Passive fiber, $f_{PE}(x)$ & Tension begins beyond optimal length. At $x=1.40$ it is 0.263, and at $x=1.70$ it reaches one. Further stretch follows the terminal slope. \\
Fiber velocity, $f_V(v)$ & At $v=-0.5$, 0, and $+0.5$, the multipliers are 0.158, 1, and 1.331. At $v=-1$ the active fiber produces zero force, and at $v=+1$ the multiplier is 1.4. \\
Tendon, $f_T(y)$ & The slack tendon carries zero tension. At 2\,\% elongation ($y=1.02$), tension is 0.261. It reaches one at 4.9\,\% elongation and then increases linearly. \\
\bottomrule
\end{tabularx}
\end{table}

At fixed total path length, series elasticity allows the tendon to lengthen as active fibers shorten. That change moves the fibers along both their force--length and force--velocity curves. Thelen and Millard give different curves for these effects. The Gait2392 analyses in this dissertation used Thelen muscles. The Millard curves provide a reference for interpreting what changes when another equilibrium formulation is used.

\subsection{OpenSim Moment Arms and Human Torque Targets}
\label{sec:app_opensim_targets}

For each Gait2392 muscle, the signed tendon-excursion moment arm about the knee is
$r_k(\theta_k)=-\partial l_{\mathrm{MT}}/\partial\theta_k$, in meters, with the knee angle in radians and the other independent coordinates fixed.
Positive tension acting through a positive moment arm produces torque in the positive knee-coordinate direction.
The reference torque in Section~\ref{sec:optimization} is $\tau_{k,\mathrm{ref}}(\theta_k)=r_k(\theta_k)F_{\mathrm{iso}}(\theta_k)$, in newton meters, where $F_{\mathrm{iso}}$ denotes the muscle's isometric tensile force at that pose.
The archived \nolinkurl{OpenSim_Vasti_MomentArm.txt} and \nolinkurl{OpenSim_Bifem_MomentArm.txt} contain the moment arms.
The corresponding \nolinkurl{OpenSim_Vasti_Results.txt} and \nolinkurl{OpenSim_Bifem_Results.txt} contain the torque curves.
The extensor search reads the vastus medialis torque column, and the flexor search reads the biceps femoris short-head torque column
(\texttt{buildKneeExtContext20mm.m}, lines 119--139, and \texttt{buildKneeFlexorContext20mm.m}, lines 129--134).
For a straight path crossing one planar revolute joint, the signed moment arm obtained from Equation~\eqref{eq:app_perpfoot} below agrees with the tendon-excursion definition.
OpenSim evaluates the human muscle paths with their moving points and couplings.

\section{Converted MuJoCo Muscle Formulation}
\label{sec:muscle-models}

The converted model represents each muscle by its routed musculotendon path but changes the calculation of tension. MuJoCo's built-in muscle assigns a constant length to the tendon and has no pennation angle. Once the path length and speed are known, the force calculation has no separate fiber-length state to solve \citep{mujoco_muscle_docs_237}.

Let $L$ denote the routed path length in meters. The conversion supplies the path-length interval $[L_-,L_+]$ and the corresponding normalized muscle lengths $[r_-,r_+]$. The effective optimal length $L_0$, constant tendon length $L_T$, normalized fiber length $x$, and normalized speed $v$ follow from
\begin{equation}
 L_0=\frac{L_+-L_-}{r_+-r_-},\qquad L_T=L_- -r_-L_0,\qquad
 x=\frac{L-L_T}{L_0},\qquad v=\frac{\dot L}{L_0 v_{\max}}.
 \label{eq:app_mj_normalize}
\end{equation}
Here $v_{\max}$ is the maximum shortening speed in optimal fiber lengths per second. MuJoCo combines an activation-dependent tensile force and a passive tensile force. Its signed actuator force, with pulling negative, is
\begin{equation}
 F_{\mathrm{act}}=-a F_A g_L(x)g_V(v)-F_P g_P(x).
 \label{eq:mj-muscle}
\end{equation}
The active and passive force scales, $F_A$ and $F_P$, are stored separately and have units of newtons. Each term uses its own saved actuator parameters \citep{mujoco_muscle_source_237}.

On the active length interval $x_{\min}\leq x\leq x_{\max}$, the normalized active force rises from zero to one at $x=1$ and returns to zero at the other end. Its quadratic segments are
\begin{equation}
 \begin{aligned}
 z&=\begin{cases}(x-x_{\min})/(1-x_{\min}),&x\leq1,\\
 (x_{\max}-x)/(x_{\max}-1),&x>1,\end{cases} &
 g_L(x)&=\begin{cases}2z^2,&z\leq\tfrac12,\\
 1-2(1-z)^2,&z>\tfrac12.\end{cases}
 \end{aligned}
 \label{eq:app_mj_length}
\end{equation}
For velocity, let $h$ be the maximum lengthening-force multiplier minus one. For passive stretch, let $b=(1+x_{\max})/2$, and let $p$ be the passive-force parameter that sets the slope and scale. The two remaining curves are
\begin{align}
 g_V(v)&=\begin{cases}
 0,&v\leq-1,\\ (1+v)^2,&-1<v\leq0,\\
 1+h-(h-v)^2/h,&0<v\leq h,\\ 1+h,&v>h,
 \end{cases}\label{eq:app_mj_velocity}\\[3pt]
 g_P(x)&=p\begin{cases}
 0,&x\leq1,\\ \tfrac12[(x-1)/(b-1)]^2,&1<x\leq b,\\
 \tfrac12+(x-b)/(b-1),&x>b.
 \end{cases}\label{eq:app_mj_passive}
\end{align}
The converted right biceps femoris short head, for example, has $x_{\min}=0$, $x_{\max}=2$, and $h=0.4$. At optimal length and zero speed it produces its active force scale at full activation. At half its maximum shortening speed the velocity multiplier falls to 0.25. In the default Millard comparison the corresponding multiplier is 0.158 (Table~\ref{tab:app_millard_curves}). Equal activation and equal normalized speed can therefore produce different tensile forces even before tendon stretch or pennation is considered.

In the walking comparison (Appendix~\ref{sec:muscle-models-compare}), the prescribed joint speeds determine the rate of change of each muscle path. The native MuJoCo and OpenSim active-force calculations use those speeds and the same activations. In a forward simulation, passive tension also contributes to the joint torque and the force available to load receptors. The pressure-dependent calculations below apply to the artificial actuators.

\section{BPA Force Model Implementation}

\subsection{Strain Measures}
The musculotendon path length $l_{LMT}$ is the sum of straight segments between route via points, transformed into a common frame at each frame crossing. The contractile strain used by the force model subtracts the artificial tendon and both end-cap fittings from the path length before normalization by the resting length,
\begin{equation}
    \varepsilon \;=\; \frac{l_{0} - \left(l_{LMT} - l_{ten} - 2\,l_{fitting} - \chi_{0}\right)}{l_{0}},
    \label{eq:app_strain}
\end{equation}
where $\chi_{0}$ is the constant length offset of Section~\ref{sec:improved_model} (zero in the baseline model). The tendon and the two end fittings occupy part of the origin-to-insertion path. Only the remaining BPA length enters the strain calculation. The constructor in \nolinkurl{MonoPamDataExplicit.m} sets each fitting to \qty{25.4}{\mm} by default. Its \texttt{get.Contraction} method and the pinned-flexor evaluator \nolinkurl{minimizeFlxPin.m} implement this subtraction.

The maximum free-load contraction at \qty{620}{\kPa} is $\varepsilon_{620} = (l_{0}-l_{620})/l_{0}$, where $l_{620}$ is the BPA length at \qty{620}{\kPa} with no external load (code variable \texttt{KMAX}; \texttt{MonoPamDataExplicit.m} line 278). The relative strain is $\varepsilon^{*} = \varepsilon / \varepsilon_{620}$ (line 281).

\subsection{Normalized Force Surface}
The normalized isometric force surface of Equation~(\ref{eq:Fstar}) is implemented in closed form in \nolinkurl{Functions/f_festo.m} (lines 1--26),
\begin{equation}
    F^{*} = c_{0}\left(e^{-c_{1}\varepsilon^{*}} - 1\right)
            + P^{*}\,e^{-c_{2}(\varepsilon^{*})^{2}},
    \qquad F^{*}\le 0 \;\Rightarrow\; 0,
    \label{eq:app_fstar}
\end{equation}
with $P^{*} = P/\qty{620}{\kPa}$. The zero-clamp $F^{*}\le 0 \Rightarrow 0$ is applied by the callers (e.g.\ \nolinkurl{Functions/festo4.m} line 35), not inside \nolinkurl{f_festo.m} itself. The coefficients were fitted with \texttt{fittype} nonlinear least squares and least-absolute-residual robustness (\nolinkurl{Testing_Data/bpaFits.m}, lines 45--60 and 115--140, lower bounds $[0,0,0]$) and are reported in Table~\ref{tab:Fstar}. The code additionally holds a \qty{40}{\mm} fit not used in this dissertation ($c_0 = 0.1224$, $c_1 = 10.47$, $c_2 = 2.023$; \nolinkurl{f_festo.m} lines 17--20).

At runtime the torque model may alternatively evaluate a spline interpolation of the manufacturer's datasheet curves (\nolinkurl{Functions/festo4.m}, backed by \texttt{FestoLookup.mat}), with force clamped to zero for $\varepsilon^{*}>1$ or $F^{*}<0$. Equation~(\ref{eq:app_fstar}) replaced this lookup in the improved model.

\subsection{Maximum Force}
The resting-length-dependent maximum force of Equation~(\ref{eq:Fmax_3D}) is implemented in \nolinkurl{Functions/maxBPAforce.m} (\qty{10}{\mm}: line 46; \qty{20}{\mm}: line 57; duplicate for \qty{20}{\mm} in \nolinkurl{Functions/Fmax20.m}) with the fixed \qty{7.5}{\mm} end-contact offset. The fitted \qty{20}{\mm} coefficients carry confidence intervals $a_1 \in [1.4745, 1.5011]$\,\si{\N\per\kPa} and $a_2 \in [0.0238, 0.0258]$\,\si{\per\kPa\per\metre} (\texttt{Fmax20.m} lines 26--27). Total force is the product $F = F^{*}(\varepsilon^{*},P^{*})\,F_{620}(l_{0})$, resolved into a vector along the force direction (\texttt{MonoPamDataExplicit.m} lines 283--287).

\section{Muscle Path Length and Moment Arm}

Path length accumulates the Euclidean norm of each segment between via points, applying the body-frame transformation whenever the path crosses between the femur and tibia frames (\texttt{MonoPamDataExplicit.m} lines 78--95, total at lines 122--133).

The vector from the joint origin to the nearest point on the line of action is
\begin{equation}
    \mathbf{r} \;=\; \mathbf{p} \;-\; \hat{\mathbf{u}}\left(\hat{\mathbf{u}}\cdot\mathbf{p}\right),
    \label{eq:app_perpfoot}
\end{equation}
where $\mathbf{p}$ runs from the joint origin to the crossing point and $\hat{\mathbf{u}}$ points along the tensile force (\texttt{MonoPamDataExplicit.m}, lines 156--168, \texttt{MonoMuscleData.m}, lines 111--123, \texttt{Mom} in \texttt{minimizeFlxPin.m}, and \nolinkurl{MonoPam_mult.m}, lines 186--201).
The force $\mathbf{f}=F\hat{\mathbf{u}}$ produces the moment vector $\mathbf{m}=\mathbf{r}\times\mathbf{f}$. Its signed knee-axis component is $M_{\mathbf{k}}=\hat{\mathbf{z}}_{\mathbf{k}}\cdot(\mathbf{r}\times\mathbf{f})$, so the signed moment arm per unit tensile force is $\hat{\mathbf{z}}_{\mathbf{k}}\cdot(\mathbf{r}\times\hat{\mathbf{u}})$. Equation~\eqref{eq:project_ma} uses the projected line of action and the corresponding in-plane force magnitude to obtain the same knee-axis torque. The flexor predictor also reports the unsigned in-plane length of $\mathbf{r}$ as \texttt{hypot(mA(:,1),mA(:,2))} (\texttt{predictKneeFlexor20mm.m}, line 217). The force routine clips negative tension to zero and discards moments for excessive stretching, $\varepsilon<-0.02$ (\texttt{MonoPamDataExplicit.m}, lines 283--287 and 307--311).

\section{Compliance Model Implementation}

Before the attachments deflect, Equation~(\ref{eq:Lm}) gives the muscle length as $l_{m} = l_{LMT} - \chi_{0} - l_{ten} - 2\,l_{fitting}$ (\texttt{minimizeFlxPin.m} lines 415--426; identically \texttt{minimizeExtX3.m} line 348). For the fitted two-bracket model, the insertion-bracket stiffness matrix and the projected scalar stiffness along the force direction are
\begin{equation}
    \mathbf{K}_{br} = \mathrm{diag}(\chi_{1},\,\chi_{2},\,\chi_{1}),
    \qquad
    k_{\hat{u},br} = \left(\hat{u}^{\,T}\mathbf{K}_{br}^{-1}\hat{u}\right)^{-1},
    \label{eq:app_kbr}
\end{equation}
matching Equation~\ref{eq:bktstiffness}. The origin-side bracket adds a scalar compliance $c_{b2}$ in series. The adopted pinned-flexor fit, row 77 of \nolinkurl{minimizeFlxPin10_results.mat}, used \nolinkurl{minimizeFlxPin2brk.m} with both brackets enabled and the \texttt{2trans} setting. Each bracket frame therefore used both rotations in Equation~\eqref{eq:app_bracketframe}. The stiffness arrays were $[X_1,X_2,X_1]$ at insertion and $[X_1,X_1,X_2]$ at origin (Table~\ref{tab:app_stiffness_arrays}). The separate \nolinkurl{minimizeFlxPin.m} uses only the first, $z$-axis rotation at both attachments.

The artificial tendon is modeled as an axial spring with stiffness
\begin{equation}
    k_{ten} \;=\; \frac{n\,A_{c}\,E}{L_{ten}},
    \label{eq:app_tendon}
\end{equation}
where $n$ is the number of parallel cable segments represented in the stiffness calculation. The model uses $n=2$ for the \qty{10}{\mm} actuator and $n=6$ for the \qty{20}{\mm} actuator. Each 19-strand cable has $A_c=\qty{1.51e-6}{\metre\squared}$ and $E=\qty{193e9}{\pascal}$ (\nolinkurl{MonoPam_mult.m}, function \texttt{Spr()}, lines 628--637). The evaluators obtain the spring rate through \nolinkurl{Spr(bpa_i)} (\texttt{minimizeFlxPin.m}, line 156). Series combination with the projected bracket compliance gives $k_{eq} = \left(k_{\hat{u},br}^{-1} + k_{ten}^{-1} + c_{b2}\right)^{-1}$ (\texttt{minimizeFlxPin.m}, lines 394--409).

The scalar compliance displacement $\delta$ of Equation~(\ref{eq:complianceforcebalance}) is found numerically with \texttt{fzero} on the bracket $[0,\, l_{m}-l_{620}]$,
\begin{equation}
    F_{620}(l_{0})\,F^{*}\!\left(\varepsilon^{*}(\delta),P^{*}\right) - k_{eq}\,\delta = 0,
    \label{eq:app_forcebalance}
\end{equation}
with $\delta \equiv 0$ when the relative strain before attachment deflection is at least one (\texttt{minimizeFlxPin.m} lines 426--441). The vector bracket deflection is then $\boldsymbol{\epsilon}_{br} = \mathbf{K}_{br}^{-1}\left(\norm{\mathbf{f}}\,\hat{\mathbf{u}}\right)$ (line 461), the cable stretch is $\norm{\mathbf{f}}/k_{ten}$ (line 478), and the attachment point, path, moment arm, and torque are recomputed. For the two-BPA flexor, both muscles share one force. The solver uses \texttt{fzero} when finite residuals of opposite sign bracket a root on $[0,F_{620}]$. If it cannot bracket or solve the root, it uses the mean of the two initial force estimates and bounds the result to the same interval (\nolinkurl{MonoPam_mult.m}, lines 556--584).

\section{Wrapping-Loss Implementation}

For the pinned-knee extensor, Equation~(\ref{eq:deltaL}) is implemented segment-wise over the wrap arc with radii $R_1 = \qty{12}{\mm}$ and $R_2 = \qty{40}{\mm}$ and segment boundaries $\theta_{1} = \ang{27}$ and $\theta_{2} = \ang{-23}$ (\texttt{minimizeExtX3.m}, nested \texttt{Contraction}, lines 327--370). For the biomimetic-knee \qty{20}{\mm} cell, the arc-length factor is the routed bend measure $b(\theta_{k})$ (arc length swept by the wrap, built with the \texttt{wRap} bend radius in \texttt{buildKneeFlexorRoute20mm.m} lines 265--268), giving $\Delta l = \chi_{3}\, b(\theta_{k})\,(1-\varepsilon^{*})^{2}$ (\nolinkurl{MonoPam_mult.m} lines 398--441), with a hardcoded-geometry fallback for degenerate routes (lines 428--434). The loss enters only the force-producing strain; the kinematic (prediction) strain excludes it (\texttt{minimizeExtX3.m} lines 236--240).

